% =====================================================================
\section{Порядок в коде и инструменты разработчика}
% =====================================================================
\lessonintro
  {разделять программу на файлы и классы, хранить историю изменений в Git и знать инструменты, которые понадобятся в большом проекте}
  {урок 3 (кнопка с подавлением дребезга), урок 14 (диагностика), приложение «Шпаргалка по C++»}
  {кнопка и светодиод оформлены как классы в отдельных файлах; у проекта есть история изменений}
  {компьютер с установленным Git; новых деталей не нужно}

\subsection{Сначала: программа растёт быстрее, чем кажется}

Скетч на сорок строк помещается в голове целиком. Итоговый проект --- это несколько сотен строк, три
датчика, экран и сеть. Если всё это лежит в одном файле с десятком глобальных переменных, любая правка
что-нибудь ломает, а вчерашнюю рабочую версию уже не вернуть. Этот урок не про новое «железо», а про
привычки, которые отличают законченный проект от брошенного: части программы лежат отдельно, у каждой
одна задача, а история изменений сохраняется.

\subsection{От функций к классам}

В уроке 3 кнопка состояла из функции \code{buttonPressed()} и трёх переменных, спрятанных внутри неё.
Для второй кнопки пришлось бы копировать всё целиком. Класс решает эту задачу: он объединяет данные
(номер вывода, время последнего изменения) и действия над ними (\code{begin()}, \code{pressed()})
в один новый тип. Сколько нужно кнопок --- столько создаём объектов.

\codecap{Кнопка и светодиод как классы}
\codeinput{l28_classes}

Раздел \code{private} --- это «внутренности»: к ним нельзя обратиться снаружи, а значит, нельзя и
случайно испортить. Всё общение с объектом идёт через несколько функций из раздела \code{public}.
Посмотрите на \code{loop()}: одна строка читается как фраза на человеческом языке.

\subsection{Несколько файлов}

Когда класс готов и отлажен, его выносят из основного скетча. В Arduino IDE новая вкладка создаётся
кнопкой с тремя точками справа от названия скетча (\textsf{New Tab}).

\begin{figure}[H]
\centering
\begin{tikzpicture}[font=\sffamily\scriptsize,
  f/.style={blk,fill=blkmem,minimum width=34mm,minimum height=13mm,align=left}]
  \node[f] (main) at (0,0) {\textbf{project.ino}\\ \code{\#include "Button.h"}\\ \code{setup()}, \code{loop()}};
  \node[f] (h) at (5.6,1.15) {\textbf{Button.h}\\ объявление класса:\\ «что умеет кнопка»};
  \node[f] (cpp) at (5.6,-1.15) {\textbf{Button.cpp}\\ тела функций:\\ «как она это делает»};
  \node[f,fill=blkper] (cfg) at (11.2,0) {\textbf{config.h}\\ номера выводов,\\ пороги, интервалы};
  \draw[arr] (main.east) -- (h.west);
  \draw[arr] (cpp.north) -- (h.south);
  \draw[arr] (h.east) -- (cfg.west);
  \draw[arr] (main.south) |- ++(2,-1.95) -| (cfg.south);
  \node[below] at (5.6,-2.75) {стрелка означает «подключает через \code{\#include}»};
\end{tikzpicture}
\caption{Типичная структура проекта: основной файл остаётся коротким, каждый класс живёт в паре файлов,
настройки собраны в одном месте}
\end{figure}

\begin{Verbatim}[fontsize=\small,bgcolor=codebg]
// Button.h
#pragma once                  // защита от повторного подключения файла
#include <Arduino.h>

class Button {
 public:
  explicit Button(int pin);
  void begin();
  bool pressed();
 private:
  int pin_;
  bool stable_ = HIGH, last_ = HIGH;
  uint32_t changed_ = 0;
};
\end{Verbatim}

Три правила порядка, которые стоит принять сразу:
\begin{itemize}
  \item \textbf{Числа получают имена.} Номера выводов, пороги и интервалы --- константы в \code{config.h},
        а не «магические» числа, разбросанные по коду.
  \item \textbf{Одна функция --- одно дело.} Если функцию трудно назвать коротко, она делает слишком много.
  \item \textbf{Пароли не хранятся рядом с кодом.} Имя и пароль сети выносят в отдельный файл
        \code{secrets.h}, который не публикуют.
\end{itemize}

\subsection{Git: машина времени для проекта}

Git хранит снимки проекта (коммиты). К любому снимку можно вернуться, сравнить его с текущей версией
и увидеть, какая именно правка всё сломала.

\begin{figure}[H]
\centering
\begin{tikzpicture}[font=\sffamily\scriptsize,node distance=19mm]
  \node[blk,fill=blkmem,minimum width=25mm,minimum height=11mm] (w) {Рабочая папка\\ правка файлов};
  \node[blk,fill=blkmem,minimum width=25mm,minimum height=11mm,right=of w] (s) {Подготовлено\\ к снимку};
  \node[blk,fill=blkmem,minimum width=25mm,minimum height=11mm,right=of s] (l) {История\\ (коммиты)};
  \node[blk,fill=blkrf,minimum width=25mm,minimum height=11mm,right=of l] (r) {GitHub, GitVerse\\ копия в сети};
  \draw[arr] (w) -- node[above]{\code{git add}} (s);
  \draw[arr] (s) -- node[above]{\code{git commit}} (l);
  \draw[arr] (l) -- node[above]{\code{git push}} (r);
  \draw[arr] (l.south) -- ++(0,-0.7) -| node[pos=0.25,below]{\code{git restore}, \code{git diff}} (w.south);
\end{tikzpicture}
\caption{Путь изменений в Git: от правки файла до копии проекта в сети}
\end{figure}

\begin{Verbatim}[fontsize=\small,bgcolor=codebg]
git init                         # один раз: сделать папку проекта репозиторием
git add .                        # подготовить все изменения
git commit -m "Кнопка включает светодиод"    # сохранить снимок с описанием
git status                       # что изменилось с последнего снимка
git log --oneline                # список снимков
git diff                         # что именно изменено в файлах
\end{Verbatim}

Делайте коммит каждый раз, когда что-то заработало, и пишите в описании, \emph{что} теперь умеет
программа. Перед рискованной правкой коммит обязателен: вернуться к нему --- дело одной команды.

\begin{caution}
Не отправляйте в открытый репозиторий пароли Wi-Fi, ключи API и сертификаты. Добавьте \code{secrets.h}
в файл \code{.gitignore} до первого коммита: удалить пароль из уже опубликованной истории очень трудно.
\end{caution}

\subsection{Что ещё пригодится}

\begin{center}\small
\begin{tabularx}{\textwidth}{@{}l X@{}}
\toprule
Инструмент & Зачем нужен \\
\midrule
PlatformIO (расширение для VS Code) & проекты из многих файлов, автодополнение, версии библиотек записаны в файле проекта \\
Симулятор Wokwi & проверить схему и программу без платы; удобно для домашней работы \\
Логический анализатор & увидеть сигналы UART, I\textsuperscript{2}C и SPI на экране компьютера и расшифровать их \\
Встроенный USB-JTAG & пошаговая отладка: остановить программу на строке и посмотреть переменные \\
Мультиметр и осциллограф & измерить питание и форму сигнала, когда «программа верная, а не работает» \\
\bottomrule
\end{tabularx}
\end{center}

В PlatformIO настройки проекта записаны в файле \code{platformio.ini}. Достаточно скопировать папку
проекта на другой компьютер, и нужные библиотеки установятся сами:

\begin{Verbatim}[fontsize=\small,bgcolor=codebg]
[env:esp32s3]
platform = espressif32
board = esp32-s3-devkitc-1
framework = arduino
monitor_speed = 115200
lib_deps =
    adafruit/Adafruit SSD1306
    bblanchon/ArduinoJson
\end{Verbatim}

\subsection{Самостоятельный опыт: наведите порядок}

Возьмите свой самый большой скетч из прошлых уроков. Создайте для него репозиторий и сделайте первый
коммит. Затем вынесите номера выводов в \code{config.h}, а кнопку и светодиод --- в классы в отдельных
файлах. После каждого шага проверяйте, что программа работает, и делайте коммит.

\begin{trouble}
\begin{itemize}[leftmargin=*]
  \item \textbf{Ошибка \code{redefinition of class}} --- в заголовочном файле нет строки \code{\#pragma once}.
  \item \textbf{Ошибка \code{was not declared in this scope}} --- в файле \code{.cpp} забыт \code{\#include <Arduino.h>}
        или заголовок своего класса.
  \item \textbf{Git «не видит» изменения} --- файл не сохранён в редакторе или указан в \code{.gitignore}.
\end{itemize}
\end{trouble}

\begin{summary}
\begin{itemize}[leftmargin=*]
  \item Класс объединяет данные и действия; объектов можно создать сколько угодно.
  \item Большой проект делят на файлы, а настройки и пароли хранят отдельно от логики.
  \item Git сохраняет историю: коммит после каждого рабочего шага и перед каждой рискованной правкой.
\end{itemize}
\end{summary}

\begin{selfcheck}
\begin{enumerate}
  \item Чем класс удобнее набора функций и глобальных переменных?
  \item Что пишут в файле \code{.h} и что --- в файле \code{.cpp}?
  \item Когда нужно делать коммит и что писать в его описании?
\end{enumerate}
\end{selfcheck}

\begin{tasks}
\begin{enumerate}
\item \lv{1} Разбейте свой скетч на функции с понятными названиями.
\item \lv{1} Добавьте комментарии: что делает каждая функция и какие у неё параметры.
\item \lv{2} Добавьте классу \code{Led} функцию \code{blink(период)}, которая мигает без \code{delay()}.
\item \lv{2} Выделите работу с кнопкой в отдельный класс с защитой от дребезга.
\item \lv{3} Напишите класс \code{Potentiometer} с усреднением и порогом из урока 5 и вынесите его в отдельные файлы.
\item \lv{4} Перенесите один из своих скетчей в PlatformIO, добавьте \code{.gitignore} и опубликуйте проект в сети.
\end{enumerate}
\end{tasks}
